% !TeX root = ../../Book3.tex
% 纲要标题：### 22.2 典范模
% 纲要位置：计划纲要.md，第 2194 行。

\section{典范模}
\label{b3:sec:2202}

% 纲要内容：
%
% * 存在条件与定义范围
% * 局部对偶的标准完备局部版本
% * 不对任意环无条件假定典范模存在
% * 先准备剩余域的内射包络 \(E_R(\kappa)\) 和 Matlis 对偶；完备 Noether 局部环上所需的有限模与 Artin 模对偶定理在本章作为明确证明模块
% * 单个典范模的局部对偶以 \(R=S/I\)、\(S\) 为给定的 \(n\) 维完备正则局部环、\(R\) 为 \(d\) 维 Cohen–Macaulay 环为标准版本，取 \(\omega_R=\operatorname{Ext}_S^{n-d}(R,S)\)
% * 对有限 \(R\)-模写出局部上同调与 \(\operatorname{Ext}_R^{d-i}(-,\omega_R)\) 的 Matlis 对偶关系；一般非 Cohen–Macaulay 情形不能把对偶复形缩成同一典范模公式
% * 证明所需正则局部覆盖作为该版本的显式假设；一般完备局部环的覆盖存在性在附录 D 的 Cohen 结构定理中注明证明或引用，不暗用未建立的结构定理
%

\subsection{典范模的定义与存在条件}
\label{b3:subsec:2202:01}

有限维向量空间的对偶仍然有限维，取两次对偶就回到原空间。局部上同调则常常不是有限模。例如对离散赋值环 \(A=k[[t]]\)，由两项 Čech 复形得到
\[
H^1_{(t)}(A)=k((t))/k[[t]].
\]
右端含有任意高的负幂，不能由有限个元素生成。因此，要把局部上同调与有限模联系起来，首先需要一种不同于 \(\operatorname{Hom}_A(-,A)\) 的对偶。

\ProseParagraphBreak
这个对偶的取值对象将是剩余域的内射包络。它既能接受子模上同态的延拓，又能检测任意非零元素；在完备 Noether 局部环上，它把有限模与 Artin 模互相转换。典范模正是通过这一对偶，把环自身的最高次局部上同调表示成一个有限模。为了让定义中的对象都有明确含义，本节先说明问题及适用范围，在\subsecref{b3:subsec:2202:04}建立内射包络和对偶，再于\cref{b3:def:canonical-module-regular-quotient}给出正式定义。

\ProseParagraphBreak
本节要证明的存在性有一组明确的数据：给定满射 \(S\twoheadrightarrow R\)，其中 \(S\) 是 \(n\) 维完备正则局部环，\(R\) 是 \(d\) 维 Cohen–Macaulay 局部环。此时可取
\[
\omega_R=\operatorname{Ext}_S^{n-d}(R,S).
\]
右端能够用 \(R\) 的有限自由 \(S\) 消解计算，因而先给出了一个实际可构造的有限模。它为什么具有对偶性质，以及为什么余维 \(n-d\) 恰好出现，都将在证明中由局部上同调的消失区间解释。

\subsection{完备局部环上的局部对偶}
\label{b3:subsec:2202:02}

先把上述离散赋值环的例子算得更具体。令 \(A=k[[t]]\)、\(K=k((t))\)、\(E=K/A\)。第二卷第7.2节的 Baer 判据说明 \(E\) 是内射 \(A\) 模：任意非零理想为 \(t^rA\)，若 \(f(t^r)=a+A\)，则取 \(g(1)=t^{-r}a+A\)，就得到延拓 \(g:A\rightarrow E\)。零理想的映射用零映射延拓。

\ProseParagraphBreak
\(E\) 中每个非零元素都能乘以某个 \(t\) 的幂，得到非零的 \(t^{-1}\) 项。因此子模 \(kt^{-1}\cong k\) 与 \(E\) 的任意非零子模相交非零。这说明内射性没有加入与剩余域毫无关系的直和因子；这正是下文“内射包络”中包络条件的含义。

\ProseParagraphBreak
记 \(D_A(N)=\operatorname{Hom}_A(N,E)\)。这里有
\[
D_A(A)=E,
\qquad
D_A(A/t^rA)=(0:_Et^r)=t^{-r}A/A\cong A/t^rA.
\]
最后一个同构把 \(a+t^rA\) 送到 \(at^{-r}+A\)。这使有限长度模的对偶仍是有限长度模，而自由模的对偶是非有限生成的 \(E\)。本节要建立的一般关系为
\[
H^i_{\mathfrak m}(M)
\cong D_R\bigl(\operatorname{Ext}_R^{d-i}(M,\omega_R)\bigr),
\]
其中 \(R\) 满足前一小节给出的假设，\(M\) 是有限 \(R\) 模。这一式子目前只是待证明的目标；其完整陈述和证明在\cref{b3:thm:canonical-module-local-duality}。

\subsection{典范模的存在性问题}
\label{b3:subsec:2202:03}

上面的目标包含三种不同的条件。Noether 性保证有限模具有有限生成的关系，并使内射包络的元素能用极大理想的幂控制。完备性使取两次对偶后的环仍是原环；若环不完备，首先出现的是它的完备化。Cohen–Macaulay 性则把环自身的局部上同调集中在唯一的次数 \(d\)，于是一个模 \(\omega_R\) 就足以承担全部对偶公式。

\ProseParagraphBreak
这些条件不能由“局部环”三个字替代。一般非 Cohen–Macaulay 环可能同时有几个非零局部上同调模，不能把上述公式中的 \(\omega_R\) 原样用于所有次数。即使某个环具有典范模，也不表示它的全部对偶信息已经集中在该模上。\subsecref{b3:subsec:2202:06}将给出一个含有嵌入点的幂级数商环，直接检验这种失败。

\ProseParagraphBreak
此外，\(S\twoheadrightarrow R\) 是本节存在性定理的假设，不是从符号 \(R\) 自动取得的构造。Cohen 结构定理能够提供更一般的完备正则局部覆盖，但它另有证明责任；本节末尾会明确说明使用该定理时增加了什么依据。

\subsection{内射包络、Matlis 对偶与有限模和 Artin 模}
\label{b3:subsec:2202:04}

先说明如何在所有可能的内射扩张中选择一个尽量小的对象。这里使用第二卷第7.2节已证明的两个结果：每个模都可嵌入内射模；Baer 判据用理想上的同态延拓检验内射性。内射包络及以下对偶定理在本节证明。松村英之《Commutative Ring Theory》的对应论述见\cite[\S 18, Theorems 18.4--18.7, pp. 145--150]{Matsumura1989CommutativeRingTheory}。

\begin{definition}\label{b3:def:injective-hull}
设 \(R\) 为交换环，\(N\subseteq M\) 为 \(R\) 模。若 \(M\) 的每个非零子模都与 \(N\) 相交非零，则此包含为\term[benzhi-kuozhang]{本质扩张}[essential extension]。若 \(M\) 还为内射模，就称 \(N\subseteq M\) 为 \(N\) 的\term[neishe-baoluo]{内射包络}[injective hull]，记一个这样的模为 \(E_R(N)\)。
\end{definition}

本质性有一个常用的检验形式：若映射 \(f:M\rightarrow L\) 在 \(N\) 上单射，则它在整个 \(M\) 上单射，因为非零的 \(\ker f\) 必与 \(N\) 相交。

\begin{proposition}\label{b3:prop:injective-hull-existence}
设 \(R\) 为交换环，\(N\) 为 \(R\) 模。\(N\) 有内射包络，且任意两个内射包络之间存在限制在 \(N\) 上为恒等的同构。这里不声称这个同构唯一。
\end{proposition}
\begin{proof}
把 \(N\) 嵌入内射模 \(J\)。在 \(J\) 中所有本质包含 \(N\) 的子模中按包含排序；一条链的并仍本质包含 \(N\)，因为并中的非零元素已落在链的某个成员中。Zorn 引理给出极大者 \(E\)。

\(E\) 没有真本质扩张。事实上，若 \(E\subseteq X\) 本质，则本质性的传递性使 \(N\subseteq X\) 本质；内射性把 \(E\hookrightarrow J\) 延拓为 \(X\rightarrow J\)。此映射的核与 \(E\) 交零，故为零。其像在 \(J\) 中仍本质包含 \(N\)，极大性使像等于 \(E\)。映射又在 \(E\) 上为恒等，故 \(X=E\)。

再把 \(E\) 嵌入内射模 \(J'\)，选取与 \(E\) 交零的极大子模 \(L\subseteq J'\)。商模 \(J'/L\) 本质包含 \(E\)：若一个非零子模与 \(E\) 的像交零，其原像就是真包含 \(L\) 且与 \(E\) 交零的子模，矛盾。由上一段，\(J'/L=E\) 的像，故 \(J'=E\oplus L\)。内射模的直和因子内射，得存在性。

若 \(E,E'\) 都是包络，由 \(E'\) 的内射性延拓 \(N\) 上的恒等为 \(u:E\rightarrow E'\)。本质性使 \(u\) 单射，而内射子模 \(u(E)\) 在 \(E'\) 中分裂。补因子与 \(N\) 交零，本质性迫使补因子为零，所以 \(u\) 是同构。
\end{proof}

以下固定交换 Noether 局部环 \((R,\mathfrak m,k)\) 及包络 \(k\subseteq E=E_R(k)\)。

\begin{proposition}\label{b3:prop:injective-hull-local-properties}
设 \((R,\mathfrak m,k)\) 为交换 Noether 局部环，\(E=E_R(k)\)。则
\[
\operatorname{Soc}_R E=(0:_E\mathfrak m)=k,
\qquad
E=\bigcup_{r\ge1}(0:_E\mathfrak m^r).
\]
此外，对任意 \(R\) 模 \(M\) 及 \(0\ne x\in M\)，存在 \(f:M\rightarrow E\) 使 \(f(x)\ne0\)。
\end{proposition}
\begin{proof}
\((0:_E\mathfrak m)\) 是 \(k\) 向量空间；它的每条一维子空间都须与指定的 \(k\) 相交，所以它恰为该 \(k\)。

固定 \(e\in E\)。若 \(\mathfrak p\in\operatorname{Ass}_R(Re)\)，则 \(Re\) 有一个同构于 \(R/\mathfrak p\) 的子模。该子模与 \(k\) 相交非零；而 \(R/\mathfrak p\) 的每个非零元素的零化子都是 \(\mathfrak p\)，相交元素的零化子又是 \(\mathfrak m\)，故 \(\mathfrak p=\mathfrak m\)。\cref{b3:prop:minimal-primes-associated}应用于有限循环模 \(Re\)，得到 \(\sqrt{\operatorname{Ann}_R(e)}=\mathfrak m\)。\(\mathfrak m\) 有限生成，因而某个 \(\mathfrak m^r\) 包含于该零化子，得并集公式。

对非零 \(x\)，\(\operatorname{Ann}_R(x)\) 是真理想。复合
\(Rx\cong R/\operatorname{Ann}_R(x)\twoheadrightarrow k\hookrightarrow E\)
把 \(x\) 送到非零元素，再用内射性将它延拓到 \(M\)。
\end{proof}

\begin{definition}\label{b3:def:matlis-dual}
设 \((R,\mathfrak m,k)\) 为交换 Noether 局部环，固定 \(E=E_R(k)\)。对 \(R\) 模 \(M\)，模
\[
D_R(M)\coloneqq\operatorname{Hom}_R(M,E)
\]
为它的\term[Matlis-duiou]{Matlis 对偶}[Matlis dual]。同态 \(u:M\rightarrow N\) 对应预复合映射 \(D_R(u):D_R(N)\rightarrow D_R(M)\)。
\end{definition}

\(D_R\) 是反变正合函子，因为 \(E\) 内射。上一命题还说明 \(D_R(M)=0\) 蕴含 \(M=0\)。评价映射
\[
\delta_M:M\longrightarrow D_RD_R(M),
\qquad \delta_M(x)(f)=f(x)
\]
对每个模都单射。不过，满射性需要限制模的类别。

\begin{lemma}\label{b3:lem:matlis-finite-length}
设 \((R,\mathfrak m,k)\) 为交换 Noether 局部环，\(D_R\) 如上。若 \(M\) 为有限长度 \(R\) 模，则
\[
\ell_R D_R(M)=\ell_R M,
\qquad \delta_M:M\xrightarrow{\sim}D_RD_R(M).
\]
并且乘法给出环同构
\[
\widehat R\xrightarrow{\sim}\operatorname{End}_R\bigl(E_R(k)\bigr),
\qquad \widehat R=\varprojlim_r R/\mathfrak m^r.
\]
\end{lemma}
\begin{proof}
先由 \(D_R(k)=\operatorname{Soc}E\cong k\) 及正合性，对组成列长度归纳得到长度等式。评价单射的两端长度相等，故为同构。

令 \(E_r=(0:_E\mathfrak m^r)=D_R(R/\mathfrak m^r)\)。前述有限长度对偶给出
\(D_R(E_r)\cong R/\mathfrak m^r\)，且元素 \(a+\mathfrak m^r\) 对应 \(e\mapsto ae\)。由 \(E=\bigcup_rE_r\)，一个 \(E\) 上的同态等价于各 \(E_r\) 上相容的同态族。因此
\[
\operatorname{End}_R(E)
=\varprojlim_r\operatorname{Hom}_R(E_r,E)
\cong\varprojlim_rR/\mathfrak m^r.
\]
这些同构把复合变为乘法，故是环同构。
\end{proof}

完备性在这里有了具体作用：只有 \(R=\widehat R\) 时，\(D_RD_R(R)=D_R(E)\) 才由同一个环 \(R\) 表示。以下的\term[Matlis-duiou-dingli]{Matlis 对偶定理}[Matlis duality theorem]将这一等式推广到两个完整的模类别。

\begin{theorem}[Matlis 对偶]\label{b3:thm:matlis-duality}
设 \((R,\mathfrak m,k)\) 为交换完备 Noether 局部环，\(E=E_R(k)\)，\(D_R=\operatorname{Hom}_R(-,E)\)。则 \(E\) 是 Artin 模；\(D_R\) 把有限 \(R\) 模变成 Artin \(R\) 模，把 Artin \(R\) 模变成有限 \(R\) 模。对这两类中的任意模 \(M\)，评价映射 \(\delta_M\) 是同构。因此 \(D_R\) 在有限模范畴与 Artin 模范畴之间给出反变等价。
\end{theorem}
\begin{proof}
先证 \(E\) 满足降链条件。对子模 \(L\subseteq E\) 记
\[
L^\perp=\{\,a\in R\mid aL=0\,\},
\qquad
J^\perp=(0:_EJ)\quad(J\subseteq R\text{ 为理想}).
\]
显然 \(L\subseteq L^{\perp\perp}\)。若 \(e\notin L\)，元素检测性质在 \(E/L\) 上给出同态 \(f:E/L\rightarrow E\)，使 \(f(e+L)\ne0\)。其与商映射的复合是 \(E\) 的自同态，前一引理和完备性使它为乘以某个 \(a\in R\)。此时 \(aL=0\) 而 \(ae\ne0\)，所以 \(e\notin L^{\perp\perp}\)。故 \(L=L^{\perp\perp}\)。同样，对 \(a\notin J\)，在 \(R/J\) 上检测 \(a+J\) 的同态给出 \(e\in J^\perp\) 且 \(ae\ne0\)，故 \(J=J^{\perp\perp}\)。两次取垂直把 \(E\) 的子模与 \(R\) 的理想建立反向对应，\(R\) 的升链条件遂给出 \(E\) 的降链条件。

若 \(M\) 有限，由满射 \(R^r\twoheadrightarrow M\) 得单射 \(D_R(M)\hookrightarrow E^r\)，故 \(D_R(M)\) 为 Artin 模。\(R\) 的 Noether 性又给出有限展示 \(R^s\rightarrow R^r\rightarrow M\rightarrow0\)。\(D_RD_R\) 是正合协变函子，\(\delta_R\) 由前一引理是同构，所以将该展示取两次对偶并比较余核，得到 \(\delta_M\) 同构。

若 \(A\) 为 Artin 模，考察全部映射 \(A\rightarrow E^r\) 的核，其中 \(r\) 可为任意非负整数。降链条件保证这些核中有极小者 \(K\)。若 \(0\ne x\in K\)，可再加一个分量 \(A\rightarrow E\) 检测 \(x\)，把核严格缩小，矛盾。因此某个映射给出单射 \(A\hookrightarrow E^r\)。取对偶得到满射 \(R^r\twoheadrightarrow D_R(A)\)，故 \(D_R(A)\) 有限。

最后，直接计算评价可得
\[
D_R(\delta_A)\circ\delta_{D_R(A)}
=1_{D_R(A)}.
\]
因为 \(D_R(A)\) 有限，\(\delta_{D_R(A)}\) 已是同构，故 \(D_R(\delta_A)\) 同构。\(D_R\) 正合且检测非零模，遂使 \(\delta_A\) 的核和余核均为零。评价的自然性说明这些同构与一切模同态相容，给出所述反变等价。
\end{proof}

Artin 性与有限长度仍须区分。例如 \(k((t))/k[[t]]\) 是上述定理中的 Artin 模，但它包含严格上升的子模链 \(t^{-1}A/A\subsetneq t^{-2}A/A\subsetneq\cdots\)，并非有限长度模。Matlis 对偶中的两个类别恰好允许这种差别。

\subsection{Cohen–Macaulay 商环的典范模}
\label{b3:subsec:2202:05}

有了 Matlis 对偶，下一步是找出正则局部环的最高次局部上同调。为此只需极小内射消解在极大理想处的那一部分，无须先分类所有不可分解内射模。

\begin{lemma}\label{b3:lem:injective-maximal-torsion}
设 \((A,\mathfrak n,k)\) 为交换 Noether 局部环，\(E=E_A(k)\)，\(J\) 为内射 \(A\) 模。若 \(\{u_\lambda\}_{\lambda\in\Lambda}\) 为 \(\operatorname{Soc}_A J\) 的一个 \(k\) 基，则
\[
\Gamma_{\mathfrak n}(J)\cong E^{(\Lambda)}.
\]
其中 \(E^{(\Lambda)}\) 为直和，指标数等于该基座的维数。
\end{lemma}
\begin{proof}
先注意，在 Noether 环上内射模的任意直和仍内射。由 Baer 判据，只需把理想 \(L\) 上的映射延拓到 \(A\)。因为 \(L\) 有限生成，其像落在有限个直和项中；逐分量延拓后再合成即可。

把每个映射 \(k\rightarrow J\)、\(1\mapsto u_\lambda\) 延拓到 \(E\rightarrow J\)，并合成为 \(v:E^{(\Lambda)}\rightarrow J\)。直和的每个元素由有限个分量组成，且由\cref{b3:prop:injective-hull-local-properties}被某个 \(\mathfrak n\) 的幂零化。因此它的任何非零子模都含有被 \(\mathfrak n\) 零化的非零元素：取零化幂数最小的元素，再乘 \(\mathfrak n\) 的适当幂。故 \(k^{(\Lambda)}\) 在该直和中本质。\(v\) 在这一个基座上单射，遂整体单射。

其像内射，故 \(J=v(E^{(\Lambda)})\oplus J'\)。\(J'\) 的基座为零，因为 \(u_\lambda\) 已张成 \(J\) 的全部基座。若 \(\Gamma_{\mathfrak n}(J')\ne0\)，同样取最小零化幂数就会找到非零基座元素，矛盾。因此极大理想幂挠部分恰为第一直和项。
\end{proof}

\begin{proposition}\label{b3:prop:regular-local-top-cohomology}
设 \((S,\mathfrak n,k)\) 为 \(n\) 维交换 Noether 正则局部环。则
\[
H^q_{\mathfrak n}(S)=0\quad(q\ne n),
\qquad
H^n_{\mathfrak n}(S)\cong E_S(k).
\]
此处不要求 \(S\) 完备，最后的同构依赖于一个剩余域基座的识别。
\end{proposition}
\begin{proof}
取 \(\mathfrak n\) 的正则生成系 \(x_1,\ldots,x_n\)，其存在由\cref{b3:thm:regular-local-characterizations}保证。\cref{b3:thm:koszul-regular-resolution}给出 \(k\) 的 Koszul 自由消解。外积配对
\(\bigwedge^qS^n\times\bigwedge^{n-q}S^n\rightarrow\bigwedge^nS^n\cong S\)
把其对偶复形识别为反向的 Koszul 复形。具体地，令 \(\Phi_q(v)(u)\) 为 \(u\wedge v\) 的最高外积系数。对 \(u\in K_{q+1}\)、\(v\in K_{n-q}\)，等式 \(d(u\wedge v)=0\) 给
\(\delta\Phi_q(v)=(-1)^q\Phi_{q+1}(dv)\)。故以 \((-1)^{q(q-1)/2}\Phi_q\) 作第 \(q\) 项的识别，所有微分均相容。因此
\[
\operatorname{Ext}_S^q(k,S)=
\begin{cases}k,&q=n,\\0,&q\ne n.\end{cases}
\]
等价地，也可将 \(n\) 个两项复形 \(S\xrightarrow{x_i}S\) 逐个对偶，每对偶一个因子便把它的非零余核提高一个次数。

现构造极小内射消解：令 \(C^0=S\)、\(J^q=E_S(C^q)\)、\(C^{q+1}=J^q/C^q\)，微分为商映射与 \(C^{q+1}\hookrightarrow J^{q+1}\) 的复合。本质性使 \(\operatorname{Soc}J^q\subseteq C^q\)，故 \(\operatorname{Hom}_S(k,J^\bullet)\) 的全部微分为零。于是其第 \(q\) 项就是 \(\operatorname{Ext}_S^q(k,S)\)。按内射消解计算 Ext 与按自由消解计算相同：对自由消解 \(P_\bullet\rightarrow k\) 及 \(J^\bullet\)，双复形 \(\operatorname{Hom}_S(P_\bullet,J^\bullet)\) 沿内射方向的正合性给前一种计算，沿自由方向的正合性给后一种计算；每个总次数只有有限项，逐次消去正合的行或列即可比较。

上述消解可在第 \(n\) 项结束。事实上，\cref{b3:thm:regular-local-global-dimension}给 \(\operatorname{Ext}_S^{n+1}(S/L,S)=0\) 对所有理想 \(L\) 成立。沿消解维数平移得 \(\operatorname{Ext}_S^1(S/L,C^n)=0\)，Baer 判据使 \(C^n\) 内射，所以 \(J^n=C^n\)、\(C^{n+1}=0\)。

前一引理与所得基座维数表明，\(\Gamma_{\mathfrak n}(J^q)\) 在 \(q\ne n\) 时为零，在 \(q=n\) 时为一份 \(E_S(k)\)。按支集挠子函子的导出函子计算局部上同调，立即得到结论。
\end{proof}

\begin{theorem}\label{b3:thm:regular-local-duality}
设 \((S,\mathfrak n,k)\) 为 \(n\) 维交换完备 Noether 正则局部环，\(E=E_S(k)\)，\(D_S=\operatorname{Hom}_S(-,E)\)。固定前一命题中的识别后，对每个有限 \(S\) 模 \(N\) 和每个整数 \(i\)，有关于 \(N\) 自然的同构
\[
H^i_{\mathfrak n}(N)
\cong D_S\bigl(\operatorname{Ext}_S^{n-i}(N,S)\bigr),
\qquad
D_S\bigl(H^i_{\mathfrak n}(N)\bigr)
\cong\operatorname{Ext}_S^{n-i}(N,S).
\]
负次 Ext 约定为零。各 \(H^i_{\mathfrak n}(N)\) 均为 Artin 模。
\end{theorem}
\begin{proof}
令 \(C^\bullet\) 为正则生成系 \(x_1,\ldots,x_n\) 的 Čech 复形，\(P_\bullet\rightarrow N\) 为长度至多 \(n\) 的有限秩自由消解。将 \(P_j\) 放在上链次数 \(-j\)，并考察总复形 \(C^\bullet\otimes_SP_\bullet\)。总微分采用
\(d(c\otimes p)=d_Cc\otimes p+(-1)^{\deg c}c\otimes d_Pp\)，故平方为零。

每个 \(C^q\) 都平坦，所以沿 \(P\) 方向取同调只留下 \(C^q\otimes_SN\)。因此总上同调为 \(H^i_{\mathfrak n}(N)\)，这里使用\cref{b3:thm:cech-local-cohomology}。换一个方向，前一命题给 \(H^q(C^\bullet)=0\) 对 \(q\ne n\)，而 \(H^n(C^\bullet)=E\)。各 \(P_j\) 自由，故这次只留下复形 \(E\otimes_SP_\bullet\)，其链次数 \(j\) 对应总次数 \(n-j\)。两次比较可直接用有限滤过完成：按一方向次数过滤总复形，把一个圈从最外侧的分量开始处理；若该行或列正合，就减去一个边界消去该分量，有限次后只剩非零同调所在的行；边界的比较同样进行。因复形在两个方向都有界，此过程终止，也与复形映射相容。于是
\[
H^i_{\mathfrak n}(N)
\cong H_{n-i}(E\otimes_SP_\bullet).
\]
有限自由性给出逐项同构
\[
E\otimes_SP_j
\cong\operatorname{Hom}_S\bigl(\operatorname{Hom}_S(P_j,S),E\bigr),
\quad e\otimes p\longmapsto[\lambda\mapsto\lambda(p)e].
\]
此同构与微分相容。\(D_S\) 正合，故右侧复形的同调为 \(D_S(\operatorname{Ext}_S^{n-i}(N,S))\)，得到第一个同构。Ext 模有限，Matlis 对偶定理使其对偶为 Artin 模；再取一次对偶并使用有限模的评价同构，得到第二个同构。
\end{proof}

现在可以准确地命名所需的有限模。

\begin{definition}\label{b3:def:canonical-module-regular-quotient}
设 \((R,\mathfrak m,k)\) 为 \(d\) 维交换完备 Noether 局部环，\(D_R=\operatorname{Hom}_R(-,E_R(k))\)。若有限 \(R\) 模 \(\omega_R\) 满足
\[
D_R(\omega_R)\cong H^d_{\mathfrak m}(R),
\]
则称 \(\omega_R\) 为 \(R\) 的\term[dianfanmo]{典范模}[canonical module]。定义要求存在这样的有限模；单独写下符号 \(\omega_R\) 不代表存在性已得到证明。
\end{definition}

Matlis 对偶定理说明，一旦典范模存在，其同构类型由 \(D_R(H^d_{\mathfrak m}(R))\) 决定。这里不指定某一个同构，所以“典范”不表示所有实现之间有唯一的同构。

\begin{proposition}\label{b3:prop:canonical-module-quotient}
设 \((S,\mathfrak n,k)\) 为 \(n\) 维交换完备 Noether 正则局部环，\(I\subsetneq S\) 为理想，\(R=S/I\)、\(\mathfrak m=\mathfrak n/I\)、\(d=\dim R\)。则
\[
\omega_R=\operatorname{Ext}_S^{n-d}(R,S)
\]
是 \(R\) 的典范模。若 \(R\) 还是 Cohen–Macaulay 环，令 \(c=n-d\)，则
\[
\operatorname{Ext}_S^q(R,S)=0\quad(q\ne c),
\qquad
\operatorname{depth}_R\omega_R=d.
\]
不同的给定完备正则局部覆盖产生同构的典范模。
\end{proposition}
\begin{proof}
先比较两个环上的对偶对象。模
\[
E_R=\operatorname{Hom}_S(R,E_S(k))=(0:_{E_S(k)}I)
\]
作为 \(R\) 模内射，因为对每个 \(R\) 模 \(M\)，评价于 \(1\) 给出
\(\operatorname{Hom}_R(M,E_R)\cong\operatorname{Hom}_S(M,E_S(k))\)，右端是正合函子。\(E_R\) 中的 \(k\) 本质，因为它已经在 \(E_S(k)\) 中本质。因此 \(E_R\) 就是 \(E_R(k)\) 的一个实现，而 \(D_R(M)=D_S(M)\)。由同一组生成元的 Čech 复形，\(H^i_{\mathfrak m}(M)=H^i_{\mathfrak n}(M)\)。正则局部对偶应用于 \(M=R\)，便给出典范模定义要求的同构；\cref{b3:lem:ext-basic-properties}说明 Ext 有限且被第一变量的零化子 \(I\) 零化，故是有限 \(R\) 模。

现假设 \(R\) 为 CM 环。\cref{b3:thm:depth-local-cohomology}与\cref{b3:thm:local-cohomology-dimension-vanishing}使 \(H^i_{\mathfrak m}(R)\) 仅在 \(i=d\) 非零。正则局部对偶立即给出 Ext 的集中次数。

进一步，\(R\) 作为 \(S\) 模的深度也是 \(d\)，因为 \(\mathfrak n\) 中元素在 \(R\) 上的作用只由其在 \(\mathfrak m\) 中的像决定。由\cref{b3:thm:auslander-buchsbaum}，\(\operatorname{pd}_SR=c\)。取长度 \(c\) 的有限自由消解 \(F_\bullet\rightarrow R\)。对偶复形 \(\operatorname{Hom}_S(F_\bullet,S)\) 仅在最后一项有上同调 \(\omega_R\)；反向排列后，它就是 \(\omega_R\) 的长度 \(c\) 的自由消解。再对偶一次，逐项双对偶回到 \(F_\bullet\)，所以
\[
\operatorname{Ext}_S^q(\omega_R,S)=
\begin{cases}R,&q=c,\\0,&q\ne c.\end{cases}
\]
故 \(\operatorname{pd}_S\omega_R=c\)，再次应用 Auslander–Buchsbaum 公式得深度 \(n-c=d\)。\(S\) 与 \(R\) 上的深度相同，得到所述结论。不同覆盖的独立性来自典范模的上述唯一同构类型。
\end{proof}

存在性一项本身没有要求 CM 性。CM 性所增加的是 Ext 的全部其他次数消失；正是这个更强结论允许下一小节把整个对偶公式写成只含 \(\omega_R\) 的形式。

\subsection{局部上同调与 Ext 的 Matlis 对偶}
\label{b3:subsec:2202:06}

为了从环境环 \(S\) 转到商环 \(R\)，需要把两个 Ext 计算联系起来。以下换环引理仅处理上一步已经出现的“只有一个非零次数”，并不要求读者先掌握一般谱序列。

\begin{lemma}\label{b3:lem:canonical-change-of-rings}
设 \(S\) 为交换 Noether 环，\(R=S/I\)。假设 \(S\) 作为自身上的模有有界内射消解，且存在整数 \(c\ge0\) 使
\(\operatorname{Ext}_S^q(R,S)=0\) 对 \(q\ne c\) 成立。记 \(W=\operatorname{Ext}_S^c(R,S)\)。则对每个有限 \(R\) 模 \(M\) 及每个 \(p\ge0\)，有自然同构
\[
\operatorname{Ext}_R^p(M,W)
\cong\operatorname{Ext}_S^{p+c}(M,S),
\]
且 \(\operatorname{Ext}_S^q(M,S)=0\) 对 \(q<c\) 成立。
\end{lemma}
\begin{proof}
取有界内射 \(S\) 消解 \(J^\bullet\) 及有限秩自由 \(R\) 消解 \(P_\bullet\rightarrow M\)。令
\[
T^q=\operatorname{Hom}_S(R,J^q),
\qquad B^{p,q}=\operatorname{Hom}_R(P_p,T^q).
\]
每个 \(T^q\) 都是内射 \(R\) 模，理由是
\(\operatorname{Hom}_R(-,T^q)\cong\operatorname{Hom}_S(-,J^q)\)
正合。\(B\) 位于第一象限；总次数 \(p+q\) 固定时只含有限个分量，微分加上通常的交替符号使两方向反交换。

先沿 \(p\) 方向取上同调。内射性使正次消失，零次为 \(\operatorname{Hom}_R(M,T^q)\cong\operatorname{Hom}_S(M,J^q)\)，所以总上同调是 \(\operatorname{Ext}_S^*(M,S)\)。再沿 \(q\) 方向取上同调。\(P_p\) 自由，故
\[
H^q(B^{p,\bullet})
=\operatorname{Hom}_R(P_p,H^q(T^\bullet)),
\]
它只在 \(q=c\) 留下 \(\operatorname{Hom}_R(P_p,W)\)。因此总次数 \(p+c\) 的上同调就是 \(\operatorname{Ext}_R^p(M,W)\)。这里的单行比较与前一定理中相同：用次数滤过逐行提升圈、消去正合行中的分量，剩余行上的边界恰是总边界的像；每个固定总次数只经过有限步。此构造与同态诱导的消解映射相容，给出自然性。
\end{proof}

\begin{theorem}[局部对偶]\label{b3:thm:canonical-module-local-duality}
设 \((S,\mathfrak n,k)\) 为 \(n\) 维交换完备 Noether 正则局部环，\(I\subsetneq S\)，且 \((R,\mathfrak m,k)=(S/I,\mathfrak n/I,k)\) 为 \(d\) 维 Cohen–Macaulay 环。取
\[
\omega_R=\operatorname{Ext}_S^{n-d}(R,S),
\qquad D_R=\operatorname{Hom}_R(-,E_R(k)).
\]
则对每个有限 \(R\) 模 \(M\) 和每个整数 \(i\)，有关于 \(M\) 自然的同构
\begin{equation}\label{b3:eq:canonical-local-duality}
\begin{aligned}
H^i_{\mathfrak m}(M)
&\cong D_R\bigl(\operatorname{Ext}_R^{d-i}(M,\omega_R)\bigr),\\
D_R\bigl(H^i_{\mathfrak m}(M)\bigr)
&\cong\operatorname{Ext}_R^{d-i}(M,\omega_R).
\end{aligned}
\end{equation}
负次局部上同调和负次 Ext 均取零。特别地，所有 \(H^i_{\mathfrak m}(M)\) 都是 Artin 模，且 \(\operatorname{Ext}_R^p(M,\omega_R)=0\) 对 \(p>d\) 成立。
\end{theorem}
\begin{proof}
令 \(c=n-d\)。\cref{b3:prop:canonical-module-quotient}给出 \(\operatorname{Ext}_S^q(R,S)\) 只在 \(q=c\) 非零。\(S\) 的有界内射消解已在\cref{b3:prop:regular-local-top-cohomology}的证明中构造。故前一引理给出
\[
\operatorname{Ext}_S^{n-i}(M,S)
\cong\operatorname{Ext}_R^{d-i}(M,\omega_R)
\]
在 \(d-i\ge0\) 时成立；若 \(d-i<0\)，左端也由引理的低次消失为零。再用\cref{b3:thm:regular-local-duality}，以及已验证的 \(D_S(M)=D_R(M)\)、\(H^i_{\mathfrak n}(M)=H^i_{\mathfrak m}(M)\)，得到两行公式。Artin 性来自正则局部对偶；当 \(i<0\) 时，局部上同调为零，第二行给出最后的 Ext 消失。
\end{proof}

上式称为\term[jubu-duiou]{局部对偶}[local duality]。它允许用有限模的消解研究通常不有限生成的局部上同调，但没有把后者断言为有限模。例如在 \(k[[t]]\) 上，\(\omega_R\cong R\)，公式的最高次仍然是 \(H^1_{(t)}(R)\cong E_R(k)\)，与节首计算一致。

\ProseParagraphBreak
最后用\subsecref{b3:subsec:2201:03}的例子检验 CM 假设的作用。取域 \(k\) 及
\[
S=k[[x,y]],\qquad R=S/(x^2,xy),\qquad\mathfrak m=(\bar x,\bar y).
\]
此前已算得
\[
\dim R=1,
\qquad H^0_{\mathfrak m}(R)=k\bar x\ne0.
\]
\cref{b3:prop:canonical-module-quotient}仍给出典范模 \(\omega_R=\operatorname{Ext}_S^1(R,S)\)，但若把局部对偶的单模公式用于 \(M=R\)、\(i=0\)，就会得到
\[
0\ne D_R\bigl(H^0_{\mathfrak m}(R)\bigr)
\cong\operatorname{Ext}_R^1(R,\omega_R)=0,
\]
矛盾。真正保留这部分信息的是另一个非零模 \(\operatorname{Ext}_S^2(R,S)\cong D_R(H^0_{\mathfrak m}(R))\)。因此一般情形须保留具有多个上同调的对偶复形，不能只取最高次局部上同调对应的那一个模。

\subsection{正则局部覆盖与 Cohen 结构定理}
\label{b3:subsec:2202:07}

本节所有存在性及局部对偶证明都从给定满射 \(S\twoheadrightarrow R\) 开始。对显式幂级数商环，这份数据已经在题目中；例如 \(k[[x_1,\ldots,x_n]]/I\) 的环境环就是所需的完备正则局部环。混合特征时也可以从给定完备离散赋值环 \(V\) 上的幂级数环 \(V[[x_1,\ldots,x_r]]\) 开始，而不把系数环误写成域。

\ProseParagraphBreak
对一个未给出展示的交换完备 Noether 局部环，Cohen 结构定理保证存在这种完备正则局部覆盖。此处将该存在性作为外部结果说明，其一般证明不属于本节：松村英之《Commutative Ring Theory》在\cite[\S 29, Theorems 29.3--29.4, pp. 225--226]{Matsumura1989CommutativeRingTheory}先建立系数环，再把环写为其上的幂级数商；本卷附录 D 安排讨论这个结构问题。本节给定覆盖的版本不需要调用这一结构定理。

\ProseParagraphBreak
若另外采用这一外部存在性结果，则每个交换完备 Noether 局部环都有本节意义下的典范模；若该环还为 CM 环，就可使用\eqref{b3:eq:canonical-local-duality}。这一推广增加的是覆盖存在性的依据，并未删除 CM 假设。至于非完备局部环，把其完备化的典范模下降为原环上的有限模，还需要另外的条件和论证，本节不作无条件的存在性断言。
